% =====================================================================
%  CONTEXTUALIZACIÓN
%  Los objetivos se transcriben LITERAL del anteproyecto §4.1 y §4.2: son los
%  aprobados y no se reescriben.
% =====================================================================
\chapter{Contextualización}

\section{Planteamiento del problema}

Orientarse en un edificio desconocido exige construir y actualizar una
representación mental del espacio. El proceso se complica cuando el entorno tiene
visibilidad limitada, estructuras repetidas y varios puntos de decisión.
\textcite{omer2025multilayer} describen esos espacios compartidos con personas
como el escenario que más exige a un robot. Vale también para quien camina sin
conocerlos.

Las edificaciones de varios pisos añaden una dificultad propia. El cambio de
nivel rompe la continuidad espacial y obliga a integrar información vertical en
el recorrido. Es el mismo problema que \textcite{kim2024indoor} resuelven para un
robot, con un grafo topológico entre plantas.

La robótica móvil colaborativa ha avanzado en exploración, logística y cobertura
distribuida. La asignación de tareas entre agentes tiene una formalización
aceptada desde el trabajo de \textcite{gerkey2004taxonomy}. Pero la mayoría de
esos desarrollos se evaluó en entornos de una sola planta, donde la segmentación
vertical no es el problema central. Los casos de limpieza doméstica y de almacén
que documenta \textcite{nair2024collaborative} son representativos de ese
escenario plano.

Resolver el cambio de piso con un solo robot resulta caro. Exige mecanismos de
locomoción capaces de subir escaleras, o intervenir el sistema de control del
ascensor. \textcite{song2024elevator} lo hacen instalando dispositivos conectados
en la propia cabina. Las dos opciones elevan el costo del hardware y
obligan a modificar el edificio.

De ahí el planteamiento de este trabajo. El problema no es cómo hacer que un
robot suba, sino \textbf{cómo lograr que dos robots sencillos se repartan el
recorrido y se pasen al usuario sin que la asistencia se interrumpa}. La
comunicación entre agentes sustituye al hardware de transición vertical.

\section{Pregunta problema}

\begin{quote}
¿Es viable mantener la continuidad de la asistencia de orientación a una persona
que cambia de nivel en un edificio, mediante un esquema de relevo entre dos
robots móviles dedicados cada uno a una planta, sin que ninguno de ellos
atraviese la zona de transición vertical?
\end{quote}

La pregunta tiene dos mitades y las dos se responden en el capítulo~4. La primera
es de funcionamiento: si el relevo ocurre. La segunda es de medida: con qué
incertidumbre se puede afirmar.

\section{Objetivos}

\subsection{Objetivo general}

Implementar un sistema colaborativo de dos robots móviles terrestres para la guía
coordinada y eficiente de asistencia a visitantes en entornos interiores con
múltiples pisos, mediante comunicación inter-robot, asignación dinámica de tareas
y una interfaz móvil de interacción usuario-robot.

\subsection{Objetivos específicos}

\begin{enumerate}[leftmargin=*,label=\textbf{OE\arabic*.}]
  \item Modelar la arquitectura funcional del sistema colaborativo de dos robots
    móviles, definiendo requerimientos, escenarios de operación con múltiples
    pisos, estrategia de asignación de tareas y esquema de comunicación
    inter-robot.
  \item Desarrollar una plataforma robótica móvil basada en dos vehículos tipo
    DonkeyCar, integrando los módulos de locomoción, sensado, procesamiento y
    comunicación necesarios para la operación colaborativa en entornos
    interiores.
  \item Programar una interfaz móvil para interacción humano-robot (HRI) basada
    en la selección de localizaciones de interés de origen y destino.
  \item Evaluar el desempeño del sistema mediante pruebas experimentales en un
    entorno interior controlado, utilizando métricas como tiempo de respuesta,
    tiempo de asignación de robot, tasa de éxito en la entrega de asistencia y
    continuidad del servicio entre niveles.
\end{enumerate}

Estos cuatro objetivos organizan el resto del documento. El capítulo~3 dedica una
sección a cada uno, y el~4 reporta lo que se midió para cada uno.
